% Gabarit pandoc des corrigés publiés (brut et guidé). Contenu sans marque ni nom de site.
% Compilation : lualatex (latexmk -lualatex).
% Le Markdown reste la source unique : ne pas éditer ce fichier .tex à la main.
\documentclass[11pt,a4paper]{article}

% --- Moteur, polices (Latin Modern) et langue ---------------------------------
\usepackage{iftex}
\ifLuaTeX\else\errmessage{Ce document se compile avec lualatex}\fi
\usepackage{amsmath,amssymb,mathtools}
\usepackage{fontspec}
% mathrm=sym : \mathrm prend ses glyphes dans la police mathématique (ℓ, μ, Ω), comme KaTeX.
\usepackage[math-style=TeX,bold-style=TeX,mathrm=sym]{unicode-math}
% Repli pour les glyphes absents de Latin Modern (indices et exposants Unicode,
% flèches, filets des schémas en texte, symboles) : DejaVu, puis Latin Modern Math
% (ℓ, lettres mathématiques Unicode des légendes).
\directlua{
  luaotfload.add_fallback("repliTexte", {"DejaVuSerif:mode=harf;", "DejaVuSans:mode=harf;", "file:latinmodern-math.otf:mode=harf;"})
  luaotfload.add_fallback("repliMono", {"DejaVuSansMono:mode=harf;", "DejaVuSans:mode=harf;", "file:latinmodern-math.otf:mode=harf;"})
}
\setmainfont{Latin Modern Roman}[RawFeature={fallback=repliTexte}]
\setsansfont{Latin Modern Sans}[RawFeature={fallback=repliTexte}]
\setmonofont{Latin Modern Mono}[RawFeature={fallback=repliMono}]
\setmathfont{Latin Modern Math}
\usepackage[french]{babel}
\frenchsetup{SmallCapsFigTabCaptions=false}

% --- Mise en page --------------------------------------------------------------
\usepackage[a4paper,margin=2.2cm,headheight=14pt,footskip=1.2cm]{geometry}
\usepackage{microtype}
\setlength{\emergencystretch}{4em}
\tolerance=1500
\setlength{\parindent}{0pt}
\setlength{\parskip}{0.55\baselineskip plus 2pt minus 1pt}
\usepackage{xcolor}
% Gris nommés : babel-french rend « ! » actif dans le document.
\definecolor{gris2}{gray}{0.98}
\definecolor{gris3}{gray}{0.97}
\definecolor{gris5}{gray}{0.95}
\definecolor{gris35}{gray}{0.65}
\definecolor{gris40}{gray}{0.60}
\definecolor{gris45}{gray}{0.55}
\definecolor{gris50}{gray}{0.50}
\definecolor{gris55}{gray}{0.45}
\definecolor{gris60}{gray}{0.40}
\definecolor{gris65}{gray}{0.35}
\definecolor{gris70}{gray}{0.30}
\definecolor{gris80}{gray}{0.20}
\definecolor{gris85}{gray}{0.15}
\usepackage{graphicx}
\usepackage{float}
\usepackage{caption}
\captionsetup{font=small,labelformat=empty,justification=centering}
\usepackage{booktabs,longtable,array,multirow,calc}
\usepackage{siunitx}
\sisetup{locale=FR}
\usepackage[version=4]{mhchem}
\usepackage{listings}
\usepackage{fancyhdr}
\usepackage{lastpage}
\usepackage[most]{tcolorbox}
\usepackage{enumitem}
\setlist{nosep,leftmargin=*,topsep=2pt}
\usepackage[hidelinks,unicode]{hyperref}
\usepackage{bookmark}
\hypersetup{pdftitle={Agro-Véto BCPST — Anglais (LVO) 2026},pdfauthor={Rédigé par IA},pdfcreator={LaTeX},pdfsubject={Corrigé brut}}
\usepackage{xurl}

% --- Pandoc : listes, tableaux, figures ----------------------------------------
\providecommand{\passthrough}[1]{#1}
\providecommand{\tightlist}{\setlength{\itemsep}{0pt}\setlength{\parskip}{0pt}}

\makeatletter
\def\maxwidth{\ifdim\Gin@nat@width>\linewidth\linewidth\else\Gin@nat@width\fi}
\def\maxheight{\ifdim\Gin@nat@height>0.55\textheight 0.55\textheight\else\Gin@nat@height\fi}
\makeatother
\setkeys{Gin}{width=\maxwidth,height=\maxheight,keepaspectratio}
\setcounter{secnumdepth}{-\maxdimen}
\providecommand{\pandocbounded}[1]{#1}

% --- Code (listings) -----------------------------------------------------------
\lstset{
  basicstyle=\ttfamily\small,
  columns=fullflexible,
  keepspaces=true,
  breaklines=true,
  breakatwhitespace=false,
  postbreak=\mbox{\textcolor{gris50}{$\hookrightarrow$}\space},
  frame=single,
  rulecolor=\color{gris35},
  framesep=4pt,
  xleftmargin=5pt,xrightmargin=5pt,
  backgroundcolor=\color{gris3},
  keywordstyle=\bfseries,
  commentstyle=\color{gris65},
  stringstyle=\color{gris80},
  showstringspaces=false,
  tabsize=4,
  extendedchars=true,
  aboveskip=6pt,belowskip=6pt,
}

% --- Macros des corrigés (KaTeX ↔ LaTeX) ---------------------------------------
% Les corrigés sont écrits pour KaTeX : on définit ici ce que KaTeX connaît et
% que LaTeX n'a pas, sans jamais retoucher le Markdown.
\providecommand{\R}{\mathbb{R}}
\providecommand{\N}{\mathbb{N}}
\providecommand{\Z}{\mathbb{Z}}
\providecommand{\Q}{\mathbb{Q}}
\providecommand{\C}{\mathbb{C}}
\providecommand{\dd}{\mathrm{d}}
\providecommand{\llbracket}{\lBrack}
\providecommand{\rrbracket}{\rBrack}
\providecommand{\lBrace}{\lbrace}
\providecommand{\rBrace}{\rbrace}
\providecommand{\Box}{\square}
\providecommand{\Diamond}{\lozenge}
\providecommand{\varnothing}{\emptyset}
\providecommand{\argmax}{\operatorname*{arg\,max}}
\providecommand{\argmin}{\operatorname*{arg\,min}}
\renewcommand{\leq}{\leqslant}
\renewcommand{\geq}{\geqslant}

% Équation hors texte : réduite à la largeur de la ligne si elle déborde.
\newsavebox{\boiteequation}
\newcommand{\equationajustee}[1]{%
  \sbox{\boiteequation}{$\displaystyle #1$}%
  \ifdim\wd\boiteequation>\linewidth
    \[\resizebox{\linewidth}{!}{\usebox{\boiteequation}}\]%
  \else
    \[#1\]%
  \fi}

% --- Encadrés pédagogiques (sobres, lisibles en noir et blanc) -----------------
\tcbset{
  encadre/.style={
    enhanced,breakable,sharp corners,
    colback=white,colframe=gris55,
    boxrule=0.4pt,leftrule=2.2pt,
    left=6pt,right=6pt,top=3pt,bottom=3pt,
    before skip=6pt,after skip=6pt,
    fonttitle=\bfseries\small,coltitle=black,
    attach title to upper,after title={\quad},
    parbox=false,
  }
}
\newtcolorbox{indication}[1]{encadre,colframe=gris45,colback=gris2,title={Indication~#1}}
\newtcolorbox{methode}{encadre,colframe=gris70,title={Méthode}}
\newtcolorbox{piege}{encadre,colframe=gris85,boxrule=0.8pt,leftrule=2.2pt,title={Piège}}
\newtcolorbox{jury}{encadre,colframe=gris60,colback=gris5,title={Jury}}
\newtcolorbox{encadrenote}{encadre,colframe=gris40}

% --- Titres de section ---------------------------------------------------------
\usepackage{titlesec}
\titleformat{\section}{\Large\bfseries}{}{0pt}{}
\titleformat{\subsection}{\large\bfseries}{}{0pt}{}
\titleformat{\subsubsection}{\normalsize\bfseries}{}{0pt}{}
\titlespacing*{\section}{0pt}{14pt plus 4pt minus 2pt}{6pt}
\titlespacing*{\subsection}{0pt}{10pt plus 3pt minus 2pt}{4pt}
\titlespacing*{\subsubsection}{0pt}{8pt plus 2pt minus 2pt}{2pt}

% --- En-tête et pied de page ---------------------------------------------------
\pagestyle{fancy}
\fancyhf{}
\fancyhead[L]{\small\textsc{Corrigé brut} · Agro-Véto BCPST}
\fancyhead[R]{\small Anglais 2026}
\fancyfoot[L]{\footnotesize Corrigé rédigé par IA, non officiel}
\fancyfoot[R]{\footnotesize \thepage\,/\,\pageref*{LastPage}}
\renewcommand{\headrulewidth}{0.4pt}
\renewcommand{\footrulewidth}{0pt}
\fancypagestyle{plain}{\fancyhf{}\fancyfoot[R]{\footnotesize \thepage\,/\,\pageref*{LastPage}}\fancyfoot[L]{\footnotesize Corrigé rédigé par IA, non officiel}\renewcommand{\headrulewidth}{0pt}}

\begin{document}
\thispagestyle{plain}
\begin{center}
  {\small\scshape Corrigé brut}\par\smallskip
  {\LARGE\bfseries Agro-Véto BCPST — Anglais (LVO) 2026\par}\medskip
  {\large Agro-Véto · BCPST · Anglais · 2026\par}
\end{center}
\vspace{2pt}\hrule\vspace{8pt}

\emph{Corrigé rédigé par IA, non officiel. Signalez toute erreur.}

\textbf{Problématique commune} : les deux articles (\emph{Courrier picard}, 20 décembre 2025 ; \emph{The Irish Times}, 10 janvier 2026) exposent les inquiétudes des agriculteurs européens face à l'accord UE-Mercosur. La partie I en restitue la nature, en anglais et avec ses propres mots ; la partie II demande une argumentation personnelle sur l'un des deux sujets au choix, que l'on traite ici tous les deux.

Durée : 2 heures, sans document ni appareil. Barème : partie I sur 10, partie II sur 10. Longueur imposée pour chaque partie : 220 mots ± 10 \%, soit 198 à 242 mots, nombre de mots indiqué à la fin. Le candidat ne traite qu'un des deux sujets d'expression et en indique le numéro. Nombre de mots : suites de caractères séparées par des espaces (mot à trait d'union = un mot, nombre = un mot, ponctuation seule non comptée).

\hypertarget{i-compruxe9hension-de-luxe9crit}{%
\section{I --- Compréhension de l'écrit}\label{i-compruxe9hension-de-luxe9crit}}

\hypertarget{why-are-farmers-in-france-and-ireland-worried-about-the-eu-mercosur-agreement-answer-the-question-in-your-own-words.}{%
\subsection{Why are farmers in France and Ireland worried about the EU-Mercosur agreement? Answer the question in your own words.}\label{why-are-farmers-in-france-and-ireland-worried-about-the-eu-mercosur-agreement-answer-the-question-in-your-own-words.}}

\emph{Problématique : des craintes économiques, sanitaires et politiques, que les deux articles présentent l'un comme éditorial engagé, l'autre comme reportage.}

Farmers in France and Ireland fear the EU-Mercosur agreement for economic, health and political reasons, as shown by an opinion piece in the Courrier picard (20 December 2025) and a more factual report in The Irish Times (10 January 2026).

Their first concern is price competition. As the Courrier picard notes, the deal will let more South American meat and sugar into Europe. Irish breeders quoted by The Irish Times, who have just enjoyed their first decent margins in decades, fear that 99,000 tonnes of cheaper beef, mostly high-value cuts, will wipe out their fragile profits. The gap in scale is huge: Brazil has 238 million cattle, Ireland six million.

Farmers also denounce unfair competition. According to these Irish farmers, South American cattle are neither tagged nor traced and are given hormones and antibiotics banned in Europe, while controls there are ineffective. They demand equivalent standards and fear that lower-quality, possibly unsafe meat will reach European consumers.

Finally, both articles suggest that farmers feel sacrificed by their leaders. The deal mainly serves European industry, since Germany hopes to sell more cars and its chemical sector is impatient, while Mercosur countries mostly export farm produce. Despite tractor protests in France and in Brussels, the signature was only postponed in December, and the deal was approved in January although Ireland voted against it. The Courrier picard therefore turns to consumers instead.

\textbf{(229 words)}

\hypertarget{ii-expression-uxe9crite}{%
\section{II --- Expression écrite}\label{ii-expression-uxe9crite}}

\hypertarget{sujet-1-do-you-think-consumers-should-be-responsible-for-protecting-local-agriculture-through-their-purchasing-choices-why-or-why-not}{%
\subsection{Sujet 1 --- Do you think consumers should be responsible for protecting local agriculture through their purchasing choices? Why or why not?}\label{sujet-1-do-you-think-consumers-should-be-responsible-for-protecting-local-agriculture-through-their-purchasing-choices-why-or-why-not}}

\emph{Problématique : le consommateur peut aider l'agriculture locale, mais faire peser sur lui seul cette responsabilité serait injuste et inefficace.}

\textbf{Subject 1}

In recent months, French and Irish farmers have protested against the EU-Mercosur agreement, which will bring cheaper South American beef onto European shelves. Since governments seem unable to stop it, an editorial in the Courrier picard urges shoppers to ``vote'' against the deal with their wallets. Should consumers be held responsible for protecting local agriculture?

To some extent, they should, because demand shapes supply. When shoppers choose local produce, they keep money in rural areas and reward higher animal welfare and environmental standards. In the United Kingdom, the Red Tractor logo, which many supermarkets display, helps consumers identify food produced by British farms. Buying local is therefore a concrete way of supporting farmers.

However, making consumers solely responsible would be unfair. When food prices rose by almost 20 per cent in a year in Britain in 2023, many low-income families simply could not afford to pay more. Moreover, labels are often confusing, and nobody can check where every ingredient in a ready meal comes from.

Responsibility should therefore be shared. Governments can impose the same standards on imports as on European products, and public canteens can set an example: in France, the EGalim law requires school canteens to serve a minimum share of quality and organic food. Retailers, too, can give local products more visibility.

In short, consumers can help to protect local agriculture, but they cannot carry this burden alone.

\textbf{(231 words)}

\hypertarget{sujet-2-is-balancing-economic-growth-with-the-protection-of-rural-livelihoods-possible}{%
\subsection{Sujet 2 --- Is balancing economic growth with the protection of rural livelihoods possible?}\label{sujet-2-is-balancing-economic-growth-with-the-protection-of-rural-livelihoods-possible}}

\emph{Problématique : la croissance menace souvent les campagnes, mais l'équilibre devient possible quand elle est pensée de façon durable et soutenue par des politiques publiques.}

\textbf{Subject 2}

In many countries, governments build motorways, factories and dams to stimulate growth, but these projects often take land, jobs and public services away from the countryside. Is balancing economic growth with the protection of rural livelihoods possible?

At first sight, the two goals seem incompatible, as growth tends to concentrate wealth in cities and large companies. China's Three Gorges Dam, for instance, produces huge amounts of electricity but forced more than a million people to leave their homes. In Brazil, the expansion of cattle ranching and soybean farming boosts exports while destroying the forest on which indigenous communities depend.

Yet growth does not have to be achieved at the expense of rural areas. Quality labels allow farmers to earn more from their know-how: Scotch whisky and Stilton cheese are protected names, which keeps value and jobs in their regions. Renewable energy and eco-tourism can also create local income, provided that rural communities share the profits.

Such a balance, however, depends on political choices. Ireland's National Broadband Plan, which aims to connect remote homes to high-speed internet, shows that public investment can help rural people take part in the modern economy. Similarly, the European Union's Common Agricultural Policy funds rural development as well as farming.

To conclude, the balance is possible, but only if growth is designed to be sustainable and inclusive rather than driven by profitability alone.

\textbf{(227 words)}

\end{document}
